\documentclass[10pt,a4paper]{article}
\usepackage[utf8]{inputenc}
\usepackage[T1]{fontenc}
\usepackage[spanish]{babel}
\usepackage{lmodern}
\usepackage[margin=2.2cm]{geometry}
\usepackage{enumitem}
\usepackage{fancyhdr}

\pagestyle{fancy}
\fancyhf{}
\renewcommand{\headrulewidth}{0pt}
\fancyhead[R]{\small Testing y Calidad}
\fancyfoot[C]{\small LSA Salud --- Plan PT-LSASALUD-G9-01 --- \today}

\setlist[enumerate]{itemsep=0.45em, leftmargin=*}

\begin{document}

\begin{center}
  \textbf{OBSERVACIONES Y SUGERENCIAS DEL PROYECTO FINAL ``LSA SALUD''\\ CICLO 2}
\end{center}

\vspace{0.5em}

\begin{enumerate}
  \item La regresión del 21/10/2026 verificó y cerró los 19 defectos del ciclo 1
    (\mbox{DM-LSA-001} a \mbox{DM-LSA-019}). Además se encontraron, corrigieron y cerraron 9 defectos
    nuevos (\mbox{DM-LSA-020} a \mbox{DM-LSA-028}): en total, 28 de 28 defectos cerrados y ninguno abierto.
    De los 15 casos, 7 quedaron aprobados (\mbox{TC-LSA-004-01}, \mbox{TC-LSA-007-01},
    \mbox{TC-LSA-008-01}, \mbox{TC-LSA-009-01}, \mbox{TC-LSA-011-01}, \mbox{TC-LSA-012-01} y
    \mbox{TC-LSA-013-01}), 8 quedaron como No probado por tener pasos manuales pendientes y ninguno
    resultó No aprobado.
  \item Los defectos nuevos se concentran en el reconocimiento y la captura: sin persona frente a la
    cámara o con la ventana de frames incompleta se reconocía ``dolor'' (\mbox{DM-LSA-021},
    \mbox{DM-LSA-022}), la oración se cerraba en la pausa entre señas (\mbox{DM-LSA-023}), ``dolor'' se
    confundía con ``yo'' (\mbox{DM-LSA-024}), y las muestras de captura se guardaban vacías o con baja
    frecuencia de detección (\mbox{DM-LSA-020}, \mbox{DM-LSA-027}). La regresión automatizada también
    detectó que el aviso de cámara denegada volvía a ocultarse (\mbox{DM-LSA-026}, regresión de
    \mbox{DM-LSA-007}), lo que confirma el valor de repetir los casos tras cada corrección.
  \item El requisito de latencia RNF-01.1 se cumple en el uso real: la latencia percibida del 21/10/2026
    tiene p50 de 166 ms (meta de 600 ms para la ruta con plantilla) y p95 de 945 ms (meta de 1 s), y
    la traducción con modelo de lenguaje tiene p95 de 573 ms. El margen del p95 es estrecho, y la prueba
    de interfaz con predicciones simuladas midió p50 de 649 y 747 ms porque se ejecuta sin GPU, por lo
    que no representa el uso real. El tiempo de inferencia del modelo no presenta regresión entre
    versiones: con las mismas 100 ventanas, el p95 pasó de 0,71 ms a 0,68 ms ($-4{,}9$~\%).
  \item Quedan pendientes las comprobaciones manuales, que ejecutará el equipo: señas reales de
    punta a punta (\mbox{TC-LSA-001-01}), calidad lingüística de la traducción y de la voz
    (\mbox{TC-LSA-002-01}, \mbox{TC-LSA-003-01}), dictado real (\mbox{TC-LSA-005-01}), captura con una
    persona real (\mbox{TC-LSA-006-01}), aviso de reporte con usuarios (\mbox{TC-LSA-010-01}), latencia
    percibida con una persona (\mbox{TC-LSA-014-01}) y precisión con una sesión independiente,
    estación, vocabulario, acta de validación, portabilidad y mantenibilidad
    (\mbox{TC-LSA-015-01}). Con ellas se podrán cerrar esos 8 casos.
  \item Se sugiere evaluar la precisión con una sesión de grabación independiente: el 100~\% obtenido
    corresponde a muestras del propio conjunto de entrenamiento y no mide la generalización. También
    se sugiere repetir la medición de latencia en la estación definitiva y ampliar la prueba
    automatizada de ``Finalizar atención'', que hoy solo comprueba que el botón exista.
\end{enumerate}

\end{document}
